% Part: computability
% Chapter: recursive-functions
% Section: non-pr-functions

\documentclass[../../../include/open-logic-section]{subfiles}

\begin{document}

\olfileid{cmp}{rec}{npr}
\olsection{Fungsi kang Dudu Rekursif Primitif}

Fungsi rekursif primitif ora nyakup kabeh fungsi
kang sacara intuitif komputabel. Sacara intuitif,
kita bisa nggawe dhaptar kabeh fungsi rekursif
primitif unèr, yaiku $f_0$, $f_1$, $f_2$,~\dots,
kanthi cara kang ndadekake kita bisa ngitung nilai
$f_x$ kanggo input $y$ kanthi efektif. Kanthi
tembung liya, fungsi $g(x,y)$ kang ditegesi dening
\[
g(x,y) = f_x(y)
\]
iku komputabel. Nanging fungsi iki uga komputabel:
\begin{eqnarray*}
h(x) & = & g(x,x) + 1 \\
& = & f_x(x) +1.
\end{eqnarray*}
Kanggo saben fungsi rekursif primitif $f_i$,
nilai $h$ lan $f_i$ béda ing $i$. Dadi $h$
komputabel nanging dudu rekursif primitif;
pratelan kang padha uga bener kanggo $g$.
Iki wangun ``efektif'' saka argumen
diagonalisasi Cantor.

Kita uga bisa menehi tuladha kang luwih cetha
saka fungsi komputabel kang dudu rekursif
primitif. Umpamane, notasi $g^n(x)$ nglambangake
$g(g(\dots g(x)))$, kanthi $n$ panyusunan
$g$ kabeh; banjur tegese rerangken fungsi
$g_0,g_1,\dots$ kanthi
\begin{eqnarray*}
g_0(x) & = & x+1 \\
g_{n + 1}(x) & = & g_n^x(x)
\end{eqnarray*}
Kita bisa mriksa yen saben fungsi $g_n$
rekursif primitif. Saben fungsi sabanjure
tuwuh luwih cepet tinimbang kang sakdurunge;
$g_1(x)$ padha karo $2x$, $g_2(x)$ padha
karo $2^x \cdot x$, lan $g_3(x)$ tuwuhe
kurang luwih kaya tumpukan eksponen saka
$x$ angka $2$. Fungsi Ackermann--P\'eter
pancen meh padha karo fungsi $G(x) = g_x(x)$,
lan bisa dibuktekake yen fungsi iki tuwuh
luwih cepet tinimbang fungsi rekursif
primitif endi wae. Cathetan panyunting: iterasi kaping nol yaiku identitas; tuwuh luwih cepet ing kene diwaca minangka dominasi kanggo input kang cukup gedhe, dudu kanggo saben input. Gambaran iki ora dadi bukti dominasi kang pepak.

Saiki kita bali menyang prakara enumerasi
fungsi rekursif primitif. Elinga yen kita wis
menehi notasi simbolis kanggo saben fungsi
rekursif primitif; mula cukup kanggo
ngenumerasi notasi-notasi kasebut. Kita bisa
menehi wilangan asli $\#(F)$ marang saben
notasi $F$ kanthi rekursif kaya mangkene:
\begin{eqnarray*}
\#(0) & = & \langle 0 \rangle \\
\#(S) & = & \langle 1 \rangle \\
\#(\Proj{n}{i}) & = & \langle 2, n, i \rangle \\
\#(\fn{Comp}_{k,l}[H,G_0,\dots,G_{k-1}]) & = & \langle
3,k,l,\#(H),\#(G_0),\dots,\#(G_{k-1}) \rangle \\
\#(\fn{Rec}_l[G,H]) & = & \langle 4, l, \#(G), \#(H) \rangle
\end{eqnarray*}
Ing kéné kita nggunakake kasunyatan yen saben
rerangken wilangan bisa dianggep minangka
wilangan asli, nganggo kode saka bagean sadurunge.
Asile, saben kode notasi oleh sawiji wilangan asli.
Mesthi wae, sawenehing rerangken (lan mula
sawenehing wilangan) ora cocog karo notasi;
nanging kita bisa netepake $f_i$ dadi fungsi
rekursif primitif unèr kang notasiné dikode
dening $i$, yen $i$ ngode notasi kaya mangkono;
yen ora, fungsi kasebut dadi fungsi konstan $0$.
Kanthi mangkono, kita nduwé cara eksplisit
kanggo ngenumerasi fungsi rekursif primitif
unèr.

(Pancen, sawenehing fungsi, kaya fungsi
konstan nol, bakal katon luwih saka sepisan
ing dhaptar. Iki ora mung akibat saka cara
pangodean, nanging uga amarga fungsi konstan
nol nduwé luwih saka siji notasi. Mengko kita
bakal weruh yen ora ana kanonisasi komputabel kang menehi wakil unik kanggo sembarang notasi manut kesamaan fungsi; upamane,
ora ana fungsi komputabel kang bisa mutusake
apa sawiji notasi nglambangake fungsi
konstan nol utawa ora.) Cathetan koreksi sumber SOL6-F039: iki ora ateges kabeh enumerasi komputabel kudu duwe ulangan. Shih-Chao Liu, ing An Enumeration of the Primitive Recursive Functions Without Repetition (1960, kaca 400–402), mbuktekake anane enumerasi komputabel tanpa ulangan kanggo fungsi rekursif primitif unèr. Kang ora bisa diitung yaiku nemokake wakil unik mau saka sembarang notasi fungsi.

Saiki fungsi $g(x,y)$ bisa kita anggep
ditetepake dening $f_x(y)$, kanthi $f_x$
nuduhake enumerasi kang lagi wae diandharake.
Kepiye kita ngerti yen $g(x,y)$ komputabel?
Sacara intuitif iki cetha: kanggo ngitung
$g(x,y)$, luwih dhisik ``bongkar'' $x$
lan priksa apa iku notasi kanggo fungsi
unèr. Yen iya, itung nilai fungsi kasebut
kanggo input~$y$. Cathetan panyunting: yen kode ora nglambangake notasi unèr kang sah, asil evaluator nol, manut definisi enumerasi sadurunge.

\begin{tagblock}{TMs}
\begin{digress}
Mbokmenawa sampeyan wis yakin yen (kanthi
upaya tartamtu!) kita bisa nulis program
(umpamane ing Java utawa C++) kang nindakake
iki; lan saiki kita bisa nggunakake tesis
Church--Turing, kang nyatakake yen apa wae
kang sacara intuitif komputabel bisa diitung
dening mesin Turing.

Mesthi wae, cara kang luwih langsung kanggo
nuduhake yen $g(x,y)$ komputabel yaiku
ngandharake kanthi eksplisit mesin Turing
kang ngitung fungsi kasebut. Kanthi mangkono,
kita ora perlu nggunakake tesis Church--Turing
utawa intuisi. Sedhela maneh kita bakal
nduwé piranti kang cukup kanggo nuduhake
yen $g(x,y)$ komputabel kanthi nggunakake
model komputasi kang bisa \emph{ditiru}
ing mesin Turing, yaiku fungsi rekursif.
\end{digress}
\end{tagblock}

\end{document}
